\section{Discussion and open problems}
\label{sec:discussion}

CAP-II is designed to make the dynamical bridge explicit: compilation-level objects define lapse ratios operationally, while covariant dynamics requires additional closure inputs.
This separation makes the theory falsifiable in layers: some predictions are definition-level and must hold in any realization of the compilation model; others are closure-dependent and can be tested to discriminate between constitutive choices.

\subsection{Assumption sensitivity}
\label{subsec:assumption_sensitivity}

Three sensitivities are central:
\begin{itemize}
  \item \emph{Cost-to-density map.} The exponent $p$ in Assumption~\ref{ass:cost_to_density} and the potential $V$ control how routing overhead backreacts as effective stress-energy; Remark~\ref{rem:cost_to_density_motivation} records a micro-motivated congestion viewpoint and a stability criterion for treating $p$ as a universality exponent.
  \item \emph{Gauge and slicing choices.} In the minimal closure we fix a protocol-rest shift gauge (Assumption~\ref{ass:shift_closure}); beyond this, physically preferred coordinate choices may be induced by protocol flow, and should be treated as part of the model specification when comparing time-dependent data.
  \item \emph{Readout bandwidth.} Assumption~\ref{ass:kernel_readout} controls which microscopic differences are unobservable and therefore which invariances can be claimed at the effective level.
\end{itemize}

\subsection{Relation to scalar--tensor frameworks and solar-system constraints}
\label{subsec:relation_scalar_tensor_constraints}

The minimal CAP-II closure \eqref{eq:omega_action_cap2} is Einstein gravity with an additional scalar source sector, not a Brans--Dicke-type modification of the geometric coupling.
Accordingly, the leading post-Newtonian coefficients of the metric are Einsteinian in regimes where the exterior is well-approximated by a Schwarzschild-like template.
Observable departures arise when the information scalar is light enough to carry exterior hair (Fisher/JNW-type behavior) or when its stress-energy produces measurable higher-order corrections in a benchmark window.
Appendix~\ref{app:ppn_solar_system} records the corresponding PPN positioning, the standard light-deflection/Shapiro-delay observables, and the massive-vs-light scalar regimes in terms of $(p,\rho_0,\lambda_F,V)$ through the constitutive identification.

\subsection{Relation to trace/regularization mechanisms}
\label{subsec:trace_relation}

Abel-type regularization and finite-part extraction provide a canonical way to assign constants to regulated orbit sums and traces.
In a dynamical setting, the same scheme choice appears as a renormalization condition in the effective action (Section~\ref{subsec:regularization_counterterms}), and therefore must be treated as part of the model specification rather than as an after-the-fact numerical trick.

\subsection{Open problems}
\label{subsec:open_problems}

One open problem is to derive, from microscopic protocol data, a physically preferred coordinate gauge or current interpretation in genuinely time-dependent regimes, beyond the minimal protocol-rest gauge used for the benchmark fits.
A second open problem is to integrate the dynamical closure with spectral/trace-formula mechanisms in a single unified framework without conflating the observer-interface layer with the covariant field layer.


